% ===== Block B05: 原文 p.37–46 | 第 2 章 仿射李代数与可积最高权表示 | 译手 tr-ch1c =====

\subsection{仿射李代数}

% [原文 p.37]

\noindent\textbf{2.1.1 仿射李代数.}\quad 本节回顾仿射李代数的可积最高权表示的基本事实。
关于仿射李代数可积最高权表示的细节，我们请读者参阅 Kac 的书 [Kac]。

设 $\mathfrak{g}$ 是复数域 $\C$ 上的单李代数，$\mathfrak{h}$ 为其嘉当子代数。以 $\Delta$ 记
$(\mathfrak{g},\mathfrak{h})$ 的根系。我们有根空间分解
\[
\mathfrak{g}=\mathfrak{h}\oplus\sum_{\alpha\in\Delta}\mathfrak{g}_{\alpha}.
\]

设 $\mathfrak{h}^*_{\R}$ 为 $\Delta$ 在 $\R$ 上的线性张成。固定 $\mathfrak{h}^*_{\R}$ 的一个字典序，
今后不再改动。这给出根系分解 $\Delta=\Delta_+\cup\Delta_-$，即把根系分成正根与负根。令
$\mathfrak{h}^*=\mathfrak{h}^*_{\R}\otimes_{\R}\C$，即 $\Delta$ 在 $\C$ 上的线性张成。设
$(\;\,,\;\,)$ 是单李代数 $\mathfrak{g}$ 的嘉当--Killing 型的一个常数倍。对每个元素
$\lambda\in\mathfrak{h}$，存在唯一的元素 $H_{\lambda}\in\mathfrak{h}^*$，使得
\[
\lambda(H)=(H_{\lambda},H)
\]
对所有 $H\in\mathfrak{h}$ 成立。对 $\alpha\in\Delta$，$H_{\alpha}$ 称为相应于根 $\alpha$ 的根向量。
% [Tutor Note: 原文此处两处记号颠倒："$\lambda\in\mathfrak{h}$" 应为 $\lambda\in\mathfrak{h}^*$，
%   "$H_\lambda\in\mathfrak{h}^*$" 应为 $H_\lambda\in\mathfrak{h}$；否则内积 $(H_\lambda,H)$ 无意义
%   （内积定义在 $\mathfrak{h}$ 上）。核对方法：把 (2.1) 代入即可看出 $(H_\lambda,H_\mu)$ 必须落在
%   $\mathfrak{h}$ 上。又 $H_\alpha$ 按 Kac 的记号是**对偶根**（coroot），而"根向量"专指张成
%   $\mathfrak{g}_\alpha$ 的 $E_\alpha$，原文此句把两者混同，照排并在报告中记作 ERRATA。]

在 $\mathfrak{h}^*$ 上我们由
\begin{equation}
(\lambda,\mu)=(H_{\lambda},H_{\mu}) \tag{2.1}
\end{equation}
引入一个内积。

设 $\theta$ 为最高（或最长）根，即
\[
\theta\in\Delta_+,\qquad \theta+\alpha_i\notin\Delta
\]
对所有单根 $\alpha_i$ 成立。例如，$A_{\ell}$ 型单李代数的最高根为
\[
\theta=\sum_{i=1}^{\ell}\alpha_i,
\]
而 $D_{\ell}$ 型单李代数的最高根为
\[
\alpha_1+2\sum_{i=2}^{\ell-2}\alpha_i+\alpha_{\ell-1}+\alpha_{\ell}.
\]

我们规范化内积 $(\;\,,\;\,)$ 并称之为规范化嘉当--Killing 型。规范化嘉当--Killing 型由条件
\begin{equation}
(\theta,\theta)=2 \tag{2.2}
\end{equation}
定义。
% [Tutor Note: (2.2) 把内积的标度钉死：对 $A_\ell,D_\ell,E_6,E_7,E_8$ 的长根都有
%   $(\alpha,\alpha)=2$；(2.2) 正是"长根长度平方为 $2$"这一约定。核对：例 2.1 中
%   $\mathfrak{sl}(2,\C)$ 的 $H$ 满足 $\alpha(H)=2$，而 $\Tr(H^2)=1^2+(-1)^2=2$，与 (2.2) 相符。
%   这个标度是全书归一化的基准：下文对偶 Coxeter 数 $g^*$（(2.15)）、$c_v$ 与 $\kappa$ 的公式
%   全都以 $(\theta,\theta)=2$ 为前提，换标度则 $c$ 与 $\kappa$ 都按同一因子缩放。]

% [原文 p.38]

注意到嘉当--Killing 型有如下性质：
\begin{equation}
([X,Y],Z)+(Y,[X,Z])=0. \tag{2.3}
\end{equation}

\begin{example}[2.1]
单李代数 $\mathfrak{g}=\mathfrak{sl}(2,\C)$ 由迹为零的 $2\times 2$ 矩阵组成。令
\[
H=\begin{pmatrix}1&0\\0&-1\end{pmatrix},\qquad
E=\begin{pmatrix}0&1\\0&0\end{pmatrix},\qquad
F=\begin{pmatrix}0&0\\1&0\end{pmatrix}.
\]
这些矩阵满足如下交换关系：
\[
[H,E]=2E,\qquad [H,F]=-2F,\qquad [E,F]=H.
\]
这时 $\mathfrak{h}$ 由 $H$ 张成，且 $\Delta=\{\alpha,-\alpha\}$，其中
$\alpha\in\mathfrak{h}^*$ 由 $\alpha(H)=2$ 定义。$\mathfrak{g}_{\alpha}$ 由 $E$ 张成，
$\mathfrak{g}_{-\alpha}$ 由 $F$ 张成。规范化嘉当--Killing 型由
\[
(X,Y)=\Tr(XY),\qquad X,Y\in\mathfrak{g}
\]
给出。
\end{example}

以 $\C[[\xi]]$ 与 $\C((\xi))$ 分别记 $\xi$ 的形式幂级数环与 $\xi$ 的形式洛朗幂级数域。

\begin{definition}[2.2]
与 $\mathfrak{g}$ 相伴的、$\C((\xi))$ 上的仿射李代数 $\hat{\mathfrak{g}}$ 定义为
\[
\hat{\mathfrak{g}}=\mathfrak{g}\otimes\C((\xi))\oplus\C c,
\]
其中 $c$ 是 $\hat{\mathfrak{g}}$ 的中心的一个元素，李代数结构由
\begin{equation}
[X\otimes f(\xi),Y\otimes g(\xi)]=[X,Y]\otimes f(\xi)g(\xi)
+c\cdot(X,Y)\operatorname*{Res}_{\xi=0}g(\xi)\dd f(\xi) \tag{2.4}
\end{equation}
给出，其中
\[
X,Y\in\mathfrak{g},\qquad f(\xi),g(\xi)\in\C((\xi)).
\]
\end{definition}
% [Tutor Note: (2.4) 的余循环项可自检：对换 $X\otimes f$ 与 $Y\otimes g$ 时
%   $\operatorname{Res}g\,\dd f=-\operatorname{Res}f\,\dd g$（分部积分，洛朗级数留数无边界项），
%   故 $[\;,\;]$ 反对称；$\dd f=n\xi^{n-1}\dd\xi$、$\dd g=m\xi^{m-1}\dd\xi$ 代入得余循环项只在
%   $n+m=0$ 时非零，即 $\operatorname{Res}_{\xi=0}g\,\dd f=n\delta_{n+m,0}$，这正是后文 (2.11)
%   里 $\delta_{n+m,0}$ 的来源。]

令
\begin{equation}
\hat{\mathfrak{g}}_+=\mathfrak{g}\otimes\C[[\xi]]\xi,\qquad
\hat{\mathfrak{g}}_-=\mathfrak{g}\otimes\C[\xi^{-1}]\xi^{-1}. \tag{2.5}
\end{equation}
我们把 $\hat{\mathfrak{g}}_+$ 与 $\hat{\mathfrak{g}}_-$ 视为 $\hat{\mathfrak{g}}$ 的李子代数。此外，下文
我们常常把 $\mathfrak{g}$ 与 $\mathfrak{g}\otimes 1$ 等同，从而可以把 $\mathfrak{g}$ 视为
$\hat{\mathfrak{g}}$ 的李子代数。我们有分解
\begin{equation}
\hat{\mathfrak{g}}=\hat{\mathfrak{g}}_+\oplus\mathfrak{g}\oplus\C c\oplus\hat{\mathfrak{g}}_-. \tag{2.6}
\end{equation}

下文我们自由地使用如下记号。
\begin{align*}
X(n)&:=X\otimes\xi^{n}, & X&\in\mathfrak{g},\quad n\in\Z,\\
X&=X(0)=X\otimes 1.
\end{align*}
% [Tutor Note: 注意 (2.5) 里 $\hat{\mathfrak{g}}_\pm$ 用**形式幂级数**（$\C[[\xi]]$、$\C[\xi^{-1}]$），
%   而下面 Remark 2.3 里的 $\mathfrak{g}_\pm$ 用**多项式**（$\C[\xi]$、$\C[\xi^{-1}]$）。
%   $\hat{\mathfrak{g}}$ 是"解析"版本（允许无穷和），$\mathfrak{g}_{aff}$ 是"代数"版本；
%   第 5 章把 $\hat{\mathfrak{g}}$ 的模与共形块对应时用的是 $\hat{\mathfrak{g}}$ 这一侧，
%   读后面 OPE 与正规序公式时务必记住无穷和是允许的。]

\begin{remark}[2.3]
注意通常仿射李代数是借助 $\C[\xi,\xi^{-1}]$ 定义的。即令
\[
\mathfrak{g}_{aff}:=\mathfrak{g}\otimes\C[\xi,\xi^{-1}]\oplus\C\cdot c
\]
并取交换关系 (2.4)。令
\[
\mathfrak{g}_+=\mathfrak{g}\otimes\C[\xi]\xi.
\]
这是一个李子代数，并且我们有
\[
\mathfrak{g}_{aff}=\mathfrak{g}_+\oplus\mathfrak{g}\oplus\C\cdot c\oplus\hat{\mathfrak{g}}_-.
\]
此外，李代数 $\hat{\mathfrak{g}}_+$ 包含在 $\mathfrak{g}_+$ 的 $\xi$-进完备化中。
\end{remark}

% [原文 p.39]

\noindent\textbf{2.1.2 可积最高权表示.}\quad 让我们回顾单李代数 $\mathfrak{g}$ 的表示论。一个不可约
左 $\mathfrak{g}$-模 $V_{\lambda}$ 称为以 $\lambda$ 为最高权的最高权模，如果存在非零向量
$e\in V_{\lambda}$（称为最高权向量）使得
\[
He=\lambda(H)e,\qquad Xe=0\quad\text{对所有 }X\in\mathfrak{g}_{\alpha},\ \alpha\in\Delta_+.
\]
众所周知，有限维不可约左 $\mathfrak{g}$-模都是最高权模，并且两个不可约左 $\mathfrak{g}$-模同构当且
仅当它们有相同的最高权。权 $\lambda\in\mathfrak{h}^*_{\R}$ 称为整权，如果
\[
2(\lambda,\alpha)/(\alpha,\alpha)\in\Z
\]
对任意 $\alpha\in\Delta$ 成立。权 $\lambda\in\mathfrak{h}^*_{\R}$ 称为支配权，如果
\[
w(\lambda)\le\lambda
\]
对 $\mathfrak{g}$ 的 Weyl 群 $W$ 的任意元素 $w$ 成立。以 $P_+$ 记 $\mathfrak{g}$ 的整支配权之集。
一个权 $\lambda$ 是不可约左 $\mathfrak{g}$-模的最高权，当且仅当 $\lambda\in P_+$。

固定一个正整数 $\ell$（称为水平）并令
\[
P_{\ell}=\{\lambda\in P_+\mid 0\le(\theta,\lambda)\le\ell\}.
\]
对每个元素 $\lambda\in P_{\ell}$，我们如下定义 Verma 模 $\mathcal{M}_{\lambda}$。令
\[
\hat{\mathfrak{p}}_+:=\hat{\mathfrak{g}}_+\oplus\mathfrak{g}\oplus\C\cdot c.
\]
则 $\hat{\mathfrak{p}}_+$ 是 $\hat{\mathfrak{g}}$ 的李子代数。设 $V_{\lambda}$ 是以 $\lambda$ 为最高权的
不可约左 $\mathfrak{g}$-模。$\hat{\mathfrak{p}}_+$ 在 $V_{\lambda}$ 上的作用定义为
\[
\begin{aligned}
cv&=\ell v &&\text{对所有 }v\in V_{\lambda},\\
av&=0 &&\text{对所有 }a\in\hat{\mathfrak{g}}_+\text{ 及 }v\in V_{\lambda}.
\end{aligned}
\]
令
\begin{equation}
\mathcal{M}_{\lambda}:=U(\hat{\mathfrak{g}})\otimes_{\hat{\mathfrak{p}}_+}V_{\lambda}. \tag{2.7}
\end{equation}
则 $\mathcal{M}_{\lambda}$ 是左 $\hat{\mathfrak{g}}$-模，称为 Verma 模。Verma 模
$\mathcal{M}_{\lambda}$ 不是不可约的，它包含极大真子模 $\mathcal{J}_{\lambda}$。商模
$\mathcal{H}_{\lambda}:=\mathcal{M}_{\lambda}/\mathcal{J}_{\lambda}$ 具有如下性质。

\begin{theorem}[2.4]
对每个 $\lambda\in P_{\ell}$，左 $\hat{\mathfrak{g}}$-模 $\mathcal{H}_{\lambda}$ 是满足如下性质的
唯一左 $\hat{\mathfrak{g}}$-模（称为可积最高权 $\hat{\mathfrak{g}}$-模）。
\begin{enumerate}
\item $V_{\lambda}=\{|v\rangle\in\mathcal{H}_{\lambda}\mid
\hat{\mathfrak{g}}_+|v\rangle=0\}$ 是以 $\lambda$ 为最高权的不可约左 $\mathfrak{g}$-模。
\item 中心元 $c$ 在 $\mathcal{H}_{\lambda}$ 上作用为 $\ell\cdot\id$。
\item $\mathcal{H}_{\lambda}$ 由 $V_{\lambda}$ 在 $\hat{\mathfrak{g}}_-$ 上生成，且只有一个关系
\begin{equation}
(X_{\theta}\otimes\xi^{-1})^{\ell-(\theta,\lambda)+1}|\lambda\rangle=0, \tag{2.8}
\end{equation}
其中 $X_{\theta}\in\mathfrak{g}$ 是相应于极大根 $\theta$ 的元素，$|\lambda\rangle\in V_{\lambda}$
是最高权向量。
\end{enumerate}
\end{theorem}
% [Tutor Note: 指数 $\ell-(\theta,\lambda)+1$ 的来历：$\lambda\in P_\ell$ 保证
%   $\ell-(\theta,\lambda)\ge0$，故指数至少为 $1$；(2.8) 说 $X_\theta\otimes\xi^{-1}$ 作用
%   $\ell-(\theta,\lambda)+1$ 次即归零，这是"可积"（$\mathfrak{sl}_2$ 三元组
%   $(X_\theta,X_{-\theta},H_\theta)$ 的有限维表示）的直接体现。核对：$\ell=1,\lambda=0$
%   给出 $(X_\theta\otimes\xi^{-1})^{2}|0\rangle=0$，即真空模里 $J^\alpha_{-1}$ 的二次幂为零，
%   与 $\mathfrak{sl}_2$ 的 $2$ 维表示一致。]

定理说明极大真子模 $\mathcal{J}_{\lambda}$ 由
\begin{equation}
\mathcal{J}_{\lambda}=U(\hat{\mathfrak{p}}_-)|J_{\lambda}\rangle \tag{2.9}
\end{equation}
给出，其中我们令
\begin{equation}
|J_{\lambda}\rangle=(X_{\theta}\otimes\xi^{-1})^{\ell-(\theta,\lambda)+1}|\lambda\rangle. \tag{2.10}
\end{equation}
% [Tutor Note: (2.9) 里的 $\hat{\mathfrak{p}}_-$ 在原文中并未定义（前面只定义了
%   $\hat{\mathfrak{p}}_+$）；按定理 2.4 第 3 条（$\mathcal{H}_\lambda$ 由 $V_\lambda$ 在
%   $\hat{\mathfrak{g}}_-$ 上生成）应读作 $\hat{\mathfrak{g}}_-$。照排原文，见报告 ISSUES。]

% [原文 p.40]

细节见 Kac [[Kac], (10.4.6)]。

类似地，我们有可积最低权右 $\hat{\mathfrak{g}}$-模 $\mathcal{H}^{\dagger}_{\lambda}$，将在下文
\S 2.2 中讨论。

下面我们简要回顾仿射李代数的可积最高权表示的通常处理。细节请读者参阅 Kac 的书 [Kac]。

一个左 $\hat{\mathfrak{g}}^{e}$-模 $\mathcal{M}$ 称为以
$\Lambda\in(\mathfrak{h}^{e})^*$ 为最高权的最高权模，如果存在元素
$|v\rangle\in\mathcal{M}$（称为最高权向量）使得对所有 $h\in\mathfrak{h}^{e}$ 有
\[
h|v\rangle=\Lambda(h)|v\rangle,\qquad \hat{\mathfrak{n}}^{e}_+|v\rangle=0,
\]
其中
\[
\hat{\mathfrak{n}}^{e}:=\sum_{\alpha\in\Delta_+}\mathfrak{g}_{\alpha},\qquad
\Delta_+:=\{m\delta+\gamma\mid m>0,\ \gamma\in\Delta(\mathfrak{g})\cup\{0\}\}\cup
\Delta(\mathfrak{g})_+,
\]
而
\[
\mathcal{M}=U(\hat{\mathfrak{g}}^{e})|v\rangle,
\]
其中 $U(\hat{\mathfrak{g}}^{e})$ 是 $\hat{\mathfrak{g}}^{e}$ 的泛包络代数。若令
\[
\Lambda:=a(\lambda)+\ell\delta
\]
并把导子 $D$ 与下一小节定义的 Virasoro 算子 $-L_0$ 等同（见下文 (2.12)），则 Verma 模
$\mathcal{M}_{\lambda}$ 与 $\mathcal{H}_{\lambda}$ 都是最高权 $\hat{\mathfrak{g}}^{e}$-模。一个最高权
$\hat{\mathfrak{g}}^{e}$-模 $\mathcal{M}$ 称为可积模，如果对 $\mathfrak{g}$ 的任意根
$\alpha\in\Delta$ 与任意整数 $n$，元素 $X_{\alpha}\otimes\xi^{n}$ 在 $\mathcal{M}$ 上局部幂零地
作用，即对任意元素 $|\Psi\rangle\in\mathcal{M}$，存在依赖于该元素 $|\Psi\rangle$ 的正整数 $m$
使得 $(X_{\alpha}\otimes\xi^{n})^{m}|\Psi\rangle=0$。已知不可约可积最高权左
$\hat{\mathfrak{g}}^{e}$-模在上述等同下同构于某个 $\mathcal{H}_{\lambda}$。
% [Tutor Note: 此处 $\Delta_+$ 被重新定义成 $\hat{\mathfrak{g}}^e$ 的正根集，而上一节
%   $\Delta_+$ 是 $\mathfrak{g}$ 的正根集；同一符号两种含义，原文靠 $\Delta(\mathfrak{g})$
%   区分。又 $a(\lambda)$ 与 $\delta$ 原文都未定义：按上下文 $a(\lambda)$ 是把 $\lambda$
%   零延拓到 $\mathfrak{h}^e=\mathfrak{h}\oplus\Cc\oplus\C D$ 上，$\ell\delta$ 用来记录中心元 $c$
%   以 $\ell$ 作用（$\delta$ 只能是 $c$ 的对偶方向，不能和通常的零根 $\delta$ 混为一谈）。
%   这是本块两处需要小心的地方之一。]

最高权可积 $\hat{\mathfrak{g}}$-模由如下定理刻画。

\begin{theorem}[2.5]
最高权 $\hat{\mathfrak{g}}$-模 $\mathcal{H}$ 可积，当且仅当对任意根向量
$X_{\alpha}\in\mathfrak{g}$ 与整数 $n$，$X_{\alpha}(n)$ 在 $\mathcal{H}$ 上局部幂零地作用，即对任意
元素 $|\varphi\rangle\in\mathcal{H}$ 存在正整数 $m$ 使得
\[
X_{\alpha}(n)^{k}|\varphi\rangle=0\qquad\text{对所有 }k\ge m.
\]
\end{theorem}

\subsection{能量-动量张量}

\noindent\textbf{2.2.1 Segal-Sugawara 形式与 Virasoro 代数.}\quad 下文我们自由地使用如下记号。
\begin{align*}
X(n)&=X\otimes\xi^{n}, & X&\in\mathfrak{g},\\
X(z)&=\sum_{n\in\Z}X(n)z^{-n-1},
\end{align*}
其中 $z$ 是一个变量。此时正规序 $\overset{\circ\circ}{\phantom{XY}}$ 定义为
\[
\overset{\circ\circ}{X(n)Y(m)}=
\begin{cases}
X(n)Y(m), & n<m,\\[2pt]
\dfrac{1}{2}\big(X(n)Y(m)+Y(m)X(n)\big), & n=m,\\[6pt]
Y(m)X(n), & n>m.
\end{cases}
\]
注意由 (2.4)，若 $n>m$ 且 $X=Y$，我们有
\begin{equation}
\overset{\circ\circ}{X(n)X(m)}=X(n)X(m)-n\delta_{n+m,0}(X,X)\cdot c. \tag{2.11}
\end{equation}
% [Tutor Note: (2.11) 的符号自检：$n>m$ 时 $\overset{\circ\circ}{X(n)X(m)}=X(m)X(n)$，
%   而由 (2.4) 算得 $[X(n),X(m)]=n\delta_{n+m,0}(X,X)c$，于是
%   $X(n)X(m)=X(m)X(n)+n\delta_{n+m,0}(X,X)c$，移项即得 (2.11) 的负号。
%   这也说明正规序的作用恰是"减去无限大的真空期望值"，正是后文 $L_n$ 良定义的关键。]

% [原文 p.41]

\begin{definition}[2.6]
水平为 $\ell$ 的能量-动量张量 $T(z)$ 定义为
\[
T(z)=\frac{1}{2(g^*+\ell)}\sum_{a=1}^{\dim\mathfrak{g}}
\overset{\circ\circ}{J^{a}(z)J^{a}(z)},
\]
其中 $\{J^{1},J^{2},\dots\}$ 是 $\mathfrak{g}$ 关于嘉当--Killing 型 $(\;\,,\;\,)$ 的标准正交基，
$g^*$ 是 $\mathfrak{g}$ 的对偶 Coxeter 数。（见下文 (2.15)。）
\end{definition}

令
\begin{equation}
L_{n}=\frac{1}{2(g^*+\ell)}\sum_{m\in\Z}\sum_{a=1}^{\dim\mathfrak{g}}
\overset{\circ\circ}{J^{a}(m)J^{a}(n-m)}. \tag{2.12}
\end{equation}
于是我们有展开
\[
T(z)=\sum_{n\in\Z}L_{n}z^{-n-2}.
\]
算子 $L_{n}$ 称为 Virasoro 算子，它作用在 $\mathcal{H}_{\lambda}$ 上。由 (2.11)，若 $n\ne0$，在
$L_{n}$ 的定义中我们不必使用正规序，即我们有
\[
L_{n}=\frac{1}{2(g^*+\ell)}\sum_{m\in\Z}\sum_{a=1}^{\dim\mathfrak{g}}
J^{a}(m)J^{a}(n-m).
\]
对 $n=0$，我们需要用正规序来定义 $L_{0}$。算子
\[
\frac{1}{2(g^*+\ell)}\sum_{m\in\Z}\sum_{a=1}^{\dim\mathfrak{g}}J^{a}(m)J^{a}(-m)
\]
不能作用在 $\mathcal{H}_{\lambda}$ 上，但
\begin{equation}
L_{0}=\frac{1}{2(g^*+\ell)}\sum_{a=1}^{\dim\mathfrak{g}}J^{a}J^{a}
+\frac{1}{(g^*+\ell)}\sum_{m=1}^{\infty}J^{a}(-m)J^{a}(m) \tag{2.13}
\end{equation}
可以作用在 $\mathcal{H}_{\lambda}$ 上。
% [Tutor Note: 为何 $n\ne0$ 不必正规序：$n\ne0$ 时对每个固定的 $n$，使 $J^a(m)J^a(n-m)$
%   两张量都"带负指标"的 $m$ 只有有限多个，求和实际有限，正规序与普通序只差一个常数
%   （可吸收进 $c$ 的项，而 $T(z)$ 的定义已把它按 $g^*+\ell$ 归一）。$n=0$ 时
%   $\sum_m J^a(m)J^a(-m)$ 是无限和，故必须正规序化；(2.13) 正是把 $m=0$ 的项
%   $\frac{1}{2(g^*+\ell)}\sum_a J^aJ^a$ 与 $m\ge1$ 的项分开写出的结果。核对：把 (2.12) 的
%   正规序按 $n=0$ 展开，$m>0$ 的项贡献 $J^a(-m)J^a(m)$，与 (2.13) 第二项系数
%   $1/(g^*+\ell)$（而非 $1/2(g^*+\ell)$）一致。]

$L_{0}$ 是单李代数 $\mathfrak{g}$ 的 Casimir 元到仿射李代数 $\hat{\mathfrak{g}}$ 的推广。让我们
简要讨论单李代数 $\mathfrak{g}$ 的 Casimir 元 $\Omega$。Casimir 元定义为
\[
\Omega=\sum_{a=1}^{\dim\mathfrak{g}}X_{a}X^{a}
\]
作为李代数 $\mathfrak{g}$ 的泛包络代数 $U(\mathfrak{g})$ 的元素，其中 $\{X_{a}\}$ 是
$\mathfrak{g}$ 的一组基，$\{X^{a}\}$ 是关于嘉当--Killing 型的对偶基。注意 $\Omega$ 与基的
选取无关。特别地我们有
\[
\Omega=\sum_{a=1}^{\dim\mathfrak{g}}J^{a}J^{a}.
\]

如下引理是众所周知的（例如见 [Hu, 22.1, p.118]、[Kac, \S 2.3 Theorem 2.6, p.23]）。

\begin{lemma}[2.7]
$\Omega$ 属于 $U(\mathfrak{g})$ 的中心。
\end{lemma}

% [原文 p.42]

\begin{lemma}[2.8]
设 $V_{\lambda}$ 是以 $\lambda$ 为最高权的最高权左 $\mathfrak{g}$-模。Casimir 元 $\Omega$ 在
$V_{\lambda}$ 上作用为
\[
\{(\lambda,\lambda)+2(\lambda,\rho)\}\cdot\id,
\]
其中
\begin{equation}
\rho=\frac{1}{2}\sum_{\alpha\in\Delta_+}\alpha. \tag{2.14}
\end{equation}
\end{lemma}

李代数 $\mathfrak{g}$ 本身通过伴随作用
\[
\ad(X)(Y)=[X,Y]
\]
成为不可约左 $\mathfrak{g}$-模。伴随表示的最高权是极大根 $\theta$。令
\begin{equation}
g^*=1+(\theta,\rho). \tag{2.15}
\end{equation}
数 $g^*$ 称为对偶 Coxeter 数。例如，对 $\mathfrak{g}=\mathfrak{sl}(n,\C)$，我们有 $g^*=n$。
% [Tutor Note: (2.15) 可用来核对：$\rho=\frac12\sum_{\alpha>0}\alpha$ 满足 $(\rho,\alpha_i)=1$，
%   故 $A_{n-1}$ 型（即 $\mathfrak{sl}(n,\C)$）时 $(\theta,\rho)=n-1$，得 $g^*=n$，与原文例子相符。
%   常用值：$A_\ell:\ell+1$；$B_\ell:2\ell-1$；$C_\ell:\ell+1$；$D_\ell:2\ell-2$；
%   $E_6:12$，$E_7:18$，$E_8:30$，$F_4:9$，$G_2:4$。注意 $g^*$ 只依赖 $\mathfrak{g}$，
%   下面的 $c_v$ 与 $L_n$ 归一化都要用它；若把 (2.15) 的 $\theta$ 换成别的正根会得到错的 $g^*$，
%   这正是"伴随表示的最高权是 $\theta$"一句的用处。]

\begin{corollary}[2.9]
\[
\sum_{a=1}^{\dim\mathfrak{g}}\ad J^{a}\ad J^{a}(X)
=\sum_{a=1}^{\dim\mathfrak{g}}[J^{a},[J^{a},X]]=2g^*X.
\]
\end{corollary}

\begin{lemma}[2.10]
对 $X\in\mathfrak{g}$ 与 $m,n\in\Z$，我们有
\[
\sum_{a=1}^{\dim\mathfrak{g}}\Big\{[X,J^{a}](m)J^{a}(n)+J^{a}(m)[X,J^{a}](n)\Big\}=0.
\]
\end{lemma}

\noindent\textit{证明.}\quad
\begin{align*}
\sum_{a=1}^{\dim\mathfrak{g}}[X,J^{a}](m)J^{a}(n)
&=\sum_{a=1}^{\dim\mathfrak{g}}\sum_{b=1}^{\dim\mathfrak{g}}
([X,J^{a}],J^{b})J^{b}(m)J^{a}(n)\\
&=-\sum_{a=1}^{\dim\mathfrak{g}}\sum_{b=1}^{\dim\mathfrak{g}}
(J^{a},[X,J^{b}])J^{b}(m)J^{a}(n)
&&\text{由 (2.1)--(2.2)}\\
&=-\sum_{b=1}^{\dim\mathfrak{g}}\sum_{a=1}^{\dim\mathfrak{g}}
J^{b}(m)(J^{a},[X,J^{b}])J^{a}(n)\\
&=-\sum_{b=1}^{\dim\mathfrak{g}}J^{b}(m)[X,J^{b}](n)\\
&=-\sum_{a=1}^{\dim\mathfrak{g}}J^{a}(m)[X,J^{a}](n).
\end{align*}
这就证明了引理。\hfill$\square$
% [Tutor Note: 原文第二步旁注"by (2.1-2)"，真正用到的是 (2.3) 的
%   不变性 $([X,J^a],J^b)+(J^a,[X,J^b])=0$，而不是 (2.1)--(2.2)。
%   另外第一步是把 $[X,J^a]$ 在标准正交基 $\{J^b\}$ 下展开：
%   $[X,J^a]=\sum_b([X,J^a],J^b)J^b$，第四步再把
%   $\sum_a(J^a,[X,J^b])J^a=[X,J^b]$ 收回——这两步都用到了 $\{J^a\}$ 关于
%   $(\;\,,\;\,)$ 的标准正交性，缺一不可。]

由直接计算可得如下重要的定理。

\begin{theorem}[2.11]
集合 $\{L_n\}$ 满足如下交换关系：
\begin{equation}
[L_n,L_m]=(n-m)L_{n+m}+\frac{c_v}{12}(n^{3}-n)\delta_{n+m,0}, \tag{2.16}
\end{equation}
其中
\begin{equation}
c_v=\frac{\ell\dim\mathfrak{g}}{g^*+\ell}. \tag{2.17}
\end{equation}
\end{theorem}
% [Tutor Note: (2.17) 是全书最重要的归一化之一。核对 $\mathfrak{sl}(2,\C)$：$\dim\mathfrak{g}=3$，
%   $g^*=2$，得 $c_v=3\ell/(\ell+2)$，与标准结果（$\mathfrak{sl}_2$ WZW 的中心荷
%   $c=3k/(k+2)$，$k$ 为水平）一致；$\ell\to\infty$ 时 $c_v\to\dim\mathfrak{g}=3$，
%   即经典极限。注意 $c_v<\dim\mathfrak{g}$ 恒成立，且 $c_v/12$ 而不是 $c_v$ 出现在 (2.16)，
%   这是 Virasoro 代数中心扩张的通常归一化（与 (2.16) 中 $(n^3-n)$ 配对）。]

% [原文 p.43]

于是集合 $\{L_n\}$ 构成中心荷为 $c_v$ 的 Virasoro 代数。此外我们有
\[
[L_n,X(m)]=-mX(n+m),\qquad \text{对 } X\in\mathfrak{g}.
\]

\begin{corollary}[2.12]
\[
[L_n,X(z)]=z^{n}\Big(z\frac{\dd}{\dd z}+n+1\Big)X(z).
\]
\end{corollary}

对 $X\in\mathfrak{g}$、$f=f(z)\in\C((z))$ 与
$\ell=\ell(z)\frac{\dd}{\dd z}\in\C((z))\frac{\dd}{\dd z}$，我们使用如下记号。
\begin{align*}
X[f]&=\operatorname*{Res}_{z=0}\big(X(z)f(z)\dd z\big),\\
T[\ell]&=\operatorname*{Res}_{z=0}\big(T(z)\ell(z)\dd z\big).
\end{align*}
特别地，我们有
\begin{equation}
L_0=T\Big[\xi\frac{\dd}{\dd\xi}\Big]. \tag{2.18}
\end{equation}
% [Tutor Note: (2.18) 用留数一算就出来：$T(z)=\sum_nL_nz^{-n-2}$，取 $\ell(z)=z$，
%   $T[z]=\operatorname{Res}_{z=0}\sum_nL_nz^{-n-1}\dd z=L_0$（只有 $n=0$ 的项留数为 $1$）。
%   这个"用留数把算子写成生成函数的积分"的写法在第 5 章会反复出现，值得记住。]

\begin{lemma}[2.13]
$X[f]$ 与 $T[\ell]$ 作用在 $\mathcal{H}_{\lambda}$ 上，并且我们有
\[
X[f]=X\otimes f(\xi),
\]
\begin{equation}
[T[\ell],X[f]]=-X[\ell(f)], \tag{2.19}
\end{equation}
\begin{equation}
[T[\ell],T[m]]=-T[[\ell,m]]+\frac{c_v}{12}\operatorname*{Res}_{z=0}\big(\ell'''m\,\dd z\big). \tag{2.20}
\end{equation}
\end{lemma}

\noindent\textit{证明.}\quad 这里我们给出第三个公式的证明。令
\[
\ell=\ell(\xi)\frac{\dd}{\dd\xi},\qquad
\ell(\xi)=\sum_{n=-n_0}^{\infty}\ell_n\xi^{n+1},
\]
\[
m=m(\xi)\frac{\dd}{\dd\xi},\qquad
m(\xi)=\sum_{n=-n_0}^{\infty}m_n\xi^{n+1}.
\]
于是我们有
\begin{align*}
[T[\ell],T[m]]&=\sum_{k,j}[\ell_jL_j,m_kL_k]\\
&=\sum_{k,j}\ell_jm_k[L_j,L_k]\\
&=\sum_{k,j}\ell_jm_k\Big\{(j-k)L_{j+k}+\frac{c_v}{12}(j^{3}-j)\delta_{j+k,0}\Big\}.
\end{align*}
另一方面，我们有
\[
[\ell,m]=\sum_{n}\Big(\sum_{j+k=n}\ell_jm_k(k-j)\Big)\xi^{n+1}\frac{\dd}{\dd\xi}.
\]
此外我们有
\[
\sum_{k,j}\ell_jm_k(j^{3}-j)\delta_{j+k,0}
=\sum_{j}\ell_jm_{-j}(j^{3}-j)
=\operatorname*{Res}_{z=0}\big(\ell'''(z)m(z)\dd z\big).
\]
这就证明了第三个公式。\hfill$\square$
% [Tutor Note: (2.20) 与 (2.16) 的一致性可自检：取
%   $\ell=z^{n+1}\frac{\dd}{\dd z}$、$m=z^{m+1}\frac{\dd}{\dd z}$（即 $\ell$ 对应 $L_n$），
%   则 $\ell'''=(n+1)n(n-1)z^{n-2}$，故
%   $\operatorname{Res}(\ell'''m\,\dd z)=(n+1)n(n-1)\delta_{n+m,0}=(n^3-n)\delta_{n+m,0}$，
%   (2.20) 右边第二项正是 $\frac{c_v}{12}(n^3-n)\delta_{n+m,0}$。这说明 (2.20) 的留数写法
%   与 (2.16) 的 $\delta$ 系数是同一件事；也解释了为何 (2.16) 里偏偏是 $(n^3-n)$ 这种组合。]

% [原文 p.44]

\begin{remark}[2.14]
在 (2.20) 的最后一个公式中我们可以使用基于如下等式的另一种表达式：
\[
\operatorname*{Res}_{z=0}\big(\ell'''(z)m(z)\dd z\big)
=-\operatorname*{Res}_{z=0}\big(m'''(z)\ell(z)\dd z\big)
=\frac{1}{2}\operatorname*{Res}_{z=0}\left(\begin{vmatrix}
\ell'(z)&m'(z)\\
\ell''(z)&m''(z)
\end{vmatrix}\dd z\right).
\]
\end{remark}
% [Tutor Note: 这个行列式恒等式可用留数的分部积分验证：对洛朗级数
%   $\operatorname{Res}(AB'\,\dd z)=-\operatorname{Res}(A'B\,\dd z)$（全导数留数为 $0$）。
%   先证 $\operatorname{Res}(\ell'''m)=-\operatorname{Res}(\ell''m')=\operatorname{Res}(\ell'm'')$；
%   再把 $\operatorname{Res}(\ell'm''-\ell''m')$ 拆开，两者由分部积分相等，故等于
%   $2\operatorname{Res}(\ell'm'')$，正好是 $\operatorname{Res}$ 的对称化（行列式展开即
%   $\ell'm''-m'\ell''$）。这个形式说明 (2.20) 的中心项只依赖 $\ell,m$ 的一、二阶导数，
%   因而是一个"局部"的 2-上循环。]

\noindent\textbf{2.2.2 最高权模上的滤过.}\quad 让我们在 $\mathcal{H}_{\lambda}$ 上定义一个滤过
$\{F_{\bullet}\}$。我们首先对非负整数 $d$ 定义 $\mathcal{H}_{\lambda}$ 的子空间
$\mathcal{H}_{\lambda}(d)$ 为
\begin{equation}
\mathcal{H}_{\lambda}(d)=\{|v\rangle\in\mathcal{H}_{\lambda}\mid
L_0|v\rangle=(d+\Delta_{\lambda})|v\rangle\}, \tag{2.21}
\end{equation}
其中
\[
\Delta_{\lambda}=\frac{(\lambda,\lambda)+2(\lambda,\rho)}{2(g^*+\ell)},\qquad
\rho=\frac{1}{2}\sum_{\alpha\in\Delta_+}\alpha.
\]
注意由 (2.12)，在 $V_{\lambda}$ 上算子 $L_0$ 作用为
\[
L_0|v\rangle=\frac{1}{2(g^*+\ell)}\sum_{a=1}^{\dim\mathfrak{g}}J^{a}J^{a}|v\rangle
\]
对 $|v\rangle\in\mathcal{H}_{\lambda}$ 成立。因此由引理 2.8，$L_0$ 在 $V_{\lambda}$ 上作为
$\Delta_{\lambda}$ 的纯量乘法作用。对正整数 $m$ 与 $|v\rangle\in V_{\lambda}$，我们有
\begin{align*}
L_0X(-m)|v\rangle&=L_0X(-m)|v\rangle+mX(-m)|v\rangle\\
&=(\Delta_{v}+m)X(-m)|v\rangle.
\end{align*}
因此 $X(-m)|v\rangle\in F_m\mathcal{H}_{\lambda}$。类似地，对正整数 $m_1,\dots,m_k$，我们有
\[
L_0X_1(m_1)\cdots X_k(-m_k)|v\rangle
=(\Delta_{\lambda}+m_1+\cdots+m_k)X_1(m_1)\cdots X_k(-m_k)|v\rangle,
\qquad |v\rangle\in V_{\lambda}.
\]
由此容易证明 $\mathcal{H}_{\lambda}(d)$ 是有限维向量空间，并且我们有
\[
\mathcal{H}_{\lambda}=\bigoplus_{d=0}^{\infty}\mathcal{H}_{\lambda}(d).
\]
对负整数 $-d$，我们令
\[
\mathcal{H}_{\lambda}(-d)=\{0\}.
\]
现在我们由
\begin{equation}
F_p\mathcal{H}_{\lambda}=\sum_{d=0}^{p}\mathcal{H}_{\lambda}(d) \tag{2.22}
\end{equation}
定义滤过 $\{F_p\mathcal{H}_{\lambda}\}$。注意上面定义的所有滤过都是递增的。
% [Tutor Note: 这一段有三处原文笔误，均已照排（见报告 ERRATA）：
%   第一，$L_0X(-m)|v\rangle$ 展开式的首项印成 $L_0X(-m)|v\rangle$，应为
%   $X(-m)L_0|v\rangle$（这正是用 $[L_0,X(-m)]=mX(-m)$）；第二，同一行
%   $\Delta_v$ 应为 $\Delta_\lambda$；第三，下一式的 $X_1(m_1)$ 应为 $X_1(-m_1)$
%   （$m_1>0$，且此处应出现升权算子）。核对方法：$L_0$ 的本征值必须增加
%   $m_1+\cdots+m_k$，若首项写成 $X_1(m_1)$ 则本征值反而减少。]

令
\begin{equation}
\mathcal{H}^{\dagger}_{\lambda}(d)=\Hom_{\C}\big(\mathcal{H}_{\lambda}(d),\C\big). \tag{2.23}
\end{equation}
则 $\mathcal{H}_{\lambda}$ 的对偶空间 $\mathcal{H}^{\dagger}_{\lambda}$ 定义为
\begin{equation}
\mathcal{H}^{\dagger}_{\lambda}=\Hom_{\C}(\mathcal{H}_{\lambda},\C)
=\prod_{d=0}^{\infty}\mathcal{H}^{\dagger}_{\lambda}(d). \tag{2.24}
\end{equation}

% [原文 p.45]

按我们的定义，$\mathcal{H}^{\dagger}_{\lambda}$ 是右 $\hat{\mathfrak{g}}$-模。递降滤过
$\{F^{p}\mathcal{H}^{\dagger}_{\lambda}\}$ 由
\begin{equation}
F^{p}\mathcal{H}^{\dagger}_{\lambda}=\prod_{d\ge p}\mathcal{H}^{\dagger}_{\lambda}(d) \tag{2.25}
\end{equation}
定义。存在一个典范的完美双线性配对
\begin{equation}
\langle\;|\;\rangle:\mathcal{H}^{\dagger}_{\lambda}\times\mathcal{H}_{\lambda}
\longrightarrow\C, \tag{2.26}
\end{equation}
它对每个 $a\in\hat{\mathfrak{g}}$ 满足如下等式：
\[
\langle u|av\rangle=\langle ua|v\rangle,\qquad
\text{对所有 }\langle u|\in\mathcal{H}^{\dagger}_{\lambda}\text{ 与 }|v\rangle\in\mathcal{H}_{\lambda}.
\]
注意滤过 $\{F_p\}$ 与 $\{F^{p}\}$ 分别在 $\mathcal{H}_{\lambda}$ 与
$\mathcal{H}^{\dagger}_{\lambda}$ 上定义了一个拓扑。关于这个拓扑，
$\mathcal{H}^{\dagger}_{\lambda}$ 是完备的，并且是可积最高权右 $\hat{\mathfrak{g}}$-模（最低权为
$\lambda$）的完备化。令
\[
V^{\dagger}_{\lambda}=\{\langle v|\in\mathcal{H}^{\dagger}_{\lambda}\mid
\langle v|\hat{\mathfrak{g}}_-=0\}.
\]
容易证明 $V^{\dagger}_{\lambda}=\mathcal{H}^{\dagger}_{\lambda}(0)$，且 $V^{\dagger}_{\lambda}$
是以 $\lambda$ 为最低权的不可约右 $\mathfrak{g}$-模。以 $\lambda$ 为最低权的可积最高权右
$\hat{\mathfrak{g}}$-模由 $V^{\dagger}_{\lambda}$ 在 $\hat{\mathfrak{g}}_+$ 上生成，且只有一个关系
\[
\langle\lambda|(X_{-\theta}\otimes\xi)^{\ell-(\theta,\lambda)+1}=0.
\]

\begin{lemma}[2.15]
\[
\begin{aligned}
X(m)\mathcal{H}_{\lambda}(d)&\subset\mathcal{H}_{\lambda}(d-m),\\
L_{m}\mathcal{H}_{\lambda}(d)&\subset\mathcal{H}_{\lambda}(d-m),\\
\mathcal{H}^{\dagger}_{\lambda}(d)X(m)&\subset\mathcal{H}^{\dagger}_{\lambda}(d+m),\\
\mathcal{H}^{\dagger}_{\lambda}(d)L_{m}&\subset\mathcal{H}^{\dagger}_{\lambda}(d+m).
\end{aligned}
\]
\end{lemma}

用类似的方式，我们可以在 Verma 模 $\mathcal{M}_{\lambda}$ 上引入一个滤过。

\begin{proposition}[2.16]
对单李代数 $\mathfrak{g}$ 的根向量 $X_{\alpha}\in\mathfrak{g}_{\alpha}$ 与任意元素
$f(\xi)\in\C((\xi))$，$X_{\alpha}\otimes f(\xi)$ 在 $\mathcal{H}_{\lambda}$ 与
$\mathcal{H}^{\dagger}_{\lambda}$ 上的作用都是局部幂零的。
\end{proposition}

现在让我们由
\[
X(n)\langle\Phi|:=-\langle\Phi|X(-n)
\]
在 $\mathcal{H}^{\dagger}_{\lambda}$ 上引入左 $\mathfrak{g}$-模结构。容易验证这确实在
$\mathcal{H}^{\dagger}_{\lambda}$ 上定义了左 $\mathfrak{g}$-模结构。

现在我们给出左 $\mathfrak{g}$-模 $\mathcal{H}^{\dagger}_{\lambda}$ 与
$\mathcal{H}_{\lambda^{\dagger}}$ 之间的关系。对元素 $\lambda\in P_{\ell}$ 令
\[
\lambda^{\dagger}=-w(\lambda),
\]
其中 $w$ 是单李代数 $\mathfrak{g}$ 的 Weyl 群的最长元素（换言之，$w(\Delta_+)=\Delta_-$）。
注意 $\lambda^{\dagger}$ 也由如下事实刻画：$-\lambda^{\dagger}$ 是 $\mathfrak{g}$-模
$V_{\lambda}$ 的最低权。例如，对 $\mathfrak{g}=\mathfrak{sl}(2,\C)$，我们有
$\lambda^{\dagger}=\lambda$。

\begin{lemma}[2.17]
存在双线性配对
\[
(\;|\;):\mathcal{H}_{\lambda}\times\mathcal{H}_{\lambda^{\dagger}}\to\C,
\]
它在相差一个常数倍的意义下唯一，使得我们有
\[
(X(n)u|v)+(u|X(-n)v)=0
\]
对任意 $X\in\mathfrak{g}$、$n\in\Z$、$|u\rangle\in\mathcal{H}_{\lambda}$、
$|v\rangle\in\mathcal{H}_{\lambda^{\dagger}}$ 成立，且若 $d\ne d'$，则
$(\;|\;)$ 在 $\mathcal{H}_{\lambda}(d)\times\mathcal{H}_{\lambda^{\dagger}}(d')$ 上为零。
\end{lemma}
% [Tutor Note: 双重对偶记号在这几页很密：$\mathcal{H}^{\dagger}_\lambda$ 的元素写成
%   $\langle u|$（左矢），$\mathcal{H}_\lambda$ 的元素写成 $|v\rangle$；配对
%   $\langle u|v\rangle$ 是 (2.26) 给出的"求值"配对，而 Lemma 2.17 的 $(\;|\;)$ 是
%   两个最高权模之间的不变量配对（contravariant form），二者不是同一个东西。
%   核对 $d\ne d'$ 时的零性：由 (2.21) 的分次，$(\;|\;)$ 保 $L_0$ 本征值，
%   故只能在同一 $d$ 上非零；这也是 (2.27) 归纳定义能进行下去的原因。]

% [原文 p.46]

\noindent\textit{证明.}\quad 由于 $V_{\lambda}\otimes V_{\lambda^{\dagger}}$ 按对角作用视为
$\mathfrak{g}$-模时只含一维平凡 $\mathfrak{g}$-模 $\C|0_{\lambda,\lambda^{\dagger}}\rangle$，
故双线性形式 $(\;|\;)\in\Hom_{\mathfrak{g}}(V_{\lambda}\otimes V_{\lambda^{\dagger}},\C)$
在相差常数倍的意义下唯一。假设我们已有具备所需性质的双线性形式
$(\;|\;)\in\Hom(F_p\mathcal{H}_{\lambda}\otimes F_p\mathcal{H}_{\lambda^{\dagger}},\C)$。
对元素
\[
X(-m)|u\rangle\in F_{p+1}\mathcal{H}_{\lambda},\qquad
|u\rangle\in F_p\mathcal{H}_{\lambda},\qquad m>0
\]
与元素 $|v\rangle\in F_{p+1}\mathcal{H}_{\lambda^{\dagger}}$，定义
\begin{equation}
(X(-m)u|v)=-(u|X(m)v). \tag{2.27}
\end{equation}
注意由于 $X(m)|v\rangle\in F_{p+1-m}\mathcal{H}_{\lambda^{\dagger}}$，右端已经定义好了。
容易证明这样就能定义出满足 Lemma 2.17 条件的双线性形式 $(\;|\;)$。\hfill$\square$

\begin{corollary}[2.18]
存在典范的左 $\hat{\mathfrak{g}}$-模同构
\[
\mathcal{H}^{\dagger}_{\lambda}\simeq\widehat{\mathcal{H}}_{\lambda^{\dagger}},
\]
其中 $\widehat{\mathcal{H}}_{\lambda^{\dagger}}$ 是 $\mathcal{H}_{\lambda^{\dagger}}$ 关于滤过
$\{F_p\}$ 的完备化。此外，若我们通过
\[
X(m)\langle u|=-\langle u|X(-m)
\]
把 $\mathcal{H}_{\lambda}$ 视为左 $\hat{\mathfrak{g}}$-模，则 $\mathcal{H}_{\lambda^{\dagger}}$
典范同构于左 $\hat{\mathfrak{g}}$-模 $\widehat{\mathcal{H}}_{\lambda^{\dagger}}$。
\end{corollary}
% [Tutor Note: 推论 2.18 的两处 $\mathcal{H}_\lambda$ 按上下文应带剑标（原文印作
%   $\mathcal{H}_\lambda$）：因为式中出现的是左矢 $\langle u|$，它属于
%   $\mathcal{H}^{\dagger}_\lambda$；结论一句也应是 $\mathcal{H}^{\dagger}_\lambda$ 与
%   $\widehat{\mathcal{H}}_{\lambda^{\dagger}}$ 同构。照排原文，见报告 ERRATA。
%   记住这个推论的内容：右模的完备化与另一最高权模的对偶同构，后面第 5 章把
%   共形块写成 $\mathcal{H}_\lambda$ 与 $\mathcal{H}_{\lambda^\dagger}$ 之间的不变量张量时要用。]

